\subsection{Weyl Eigenvalue Bounds}

Throughout this subsection, $n\in\mathbb{Z}_{>0}$.

\begin{theorem} \label{thm:weyl-eigenvalue-inequalities-hermitian-sums} \hlt{Weyl Eigenvalue Inequalities}\\
Let $A,B\in\operatorname{Mat}(n,n,\mathbb{C})$ be Hermitian.
\begin{enumerate}[label=(\roman*)]
\setlength{\itemsep}{0pt}
\item $(\forall\,i\in J_n)(\forall\,j\in J_{n-i}\cup\{0\})\,\lambda_i(A+B)\leq\lambda_{i+j}(A)+\lambda_{n-j}(B)$.
\item $(\forall\,i\in J_n)(\forall\,j\in J_i)\,\lambda_{i-j+1}(A)+\lambda_j(B)\leq\lambda_i(A+B)$.
\end{enumerate}
\begin{proof}
Fix increasingly indexed orthonormal eigenbases $(\mathbf{u}_k)$, $(\mathbf{v}_k)$, and $(\mathbf{w}_k)$ of $A$, $A+B$, and $B$ by Theorem \ref{thm:complex-spectral-theorem}, Definition \ref{def:increasing-hermitian-eigenvalue-ordering}(i).
\begin{enumerate}[label=(\roman*)]
\setlength{\itemsep}{0pt}
\item Fix $i\in J_n$ and $j\in J_{n-i}\cup\{0\}$. Set $\mathcal{U}=\operatorname{span}(\mathbf{u}_1,\ldots,\mathbf{u}_{i+j})$ and $\mathcal{V}=\operatorname{span}(\mathbf{v}_i,\ldots,\mathbf{v}_n)$.\\
Set $\mathcal{W}=\operatorname{span}(\mathbf{w}_1,\ldots,\mathbf{w}_{n-j})$. $\dim(\mathcal{U}\cap\mathcal{V})\geq j+1$ by Corollary \ref{cor:large-subspaces-intersect}(ii) with $\dim(\mathcal{U})=i+j$ and $\dim(\mathcal{V})=n-i+1$.\\
Fix a unit $\mathbf{x}\in\mathcal{U}\cap\mathcal{V}\cap\mathcal{W}$ by Corollary \ref{cor:large-subspaces-intersect}(i), normalization, the preceding bound, and $\dim(\mathcal{W})=n-j$.\\
The chain $\lambda_i(A+B)\leq\mathbf{x}^*(A+B)\mathbf{x}=\mathbf{x}^*A\mathbf{x}+\mathbf{x}^*B\mathbf{x}\leq\lambda_{i+j}(A)+\lambda_{n-j}(B)$ holds by Theorem \ref{thm:rayleigh-quotient-spectral-interval}.
\item Fix $i\in J_n$ and $j\in J_i$, and set $p\defeq n-i+1$ and $q\defeq i-j$. $p\in J_n$ and $q\in J_{n-p}\cup\{0\}$ by index arithmetic.\\
$(\forall\,C\in\{A,B,A+B\})(\forall\,k\in J_n)\,\lambda_k(-C)=-\lambda_{n-k+1}(C)$ by Theorem \ref{thm:complex-spectral-theorem} and Definition \ref{def:increasing-hermitian-eigenvalue-ordering}(i).\\
$-\lambda_i(A+B)\leq-\lambda_j(B)-\lambda_{i-j+1}(A)$ by part (i) for $(-B,-A,p,q)$ and the reversal identity.\\
Claim (ii) holds by multiplying this inequality by $-1$.
\end{enumerate}
\hlt{Structural Commentary}: Three-flag intersection establishes (i); spectral reversal and sign multiplication establish (ii).\\
\hlt{Key Insight}: $\dim(\mathcal{U})+\dim(\mathcal{V})+\dim(\mathcal{W})=2n+1$ forces a common Rayleigh vector.
\end{proof}
\end{theorem}

\begin{computationmethod} \label{meth:weyl-inequality-method} \hlt{Weyl Eigenvalue-Bound Method}\\
Given Hermitian $A,B\in\operatorname{Mat}(n,n,\mathbb{C})$ and $i\in J_n$, compute a certified interval for $\lambda_i(A+B)$ as follows:
\begin{enumerate}[label=(\roman*)]
\setlength{\itemsep}{0pt}
\item compute the eigenvalues of $A$ and $B$ in increasing order;
\item set $U_i\defeq\min\limits_{j\in J_{n-i}\cup\{0\}}\{\lambda_{i+j}(A)+\lambda_{n-j}(B)\}$;
\item set $L_i\defeq\max\limits_{j\in J_i}\{\lambda_{i-j+1}(A)+\lambda_j(B)\}$;
\item output $[L_i,U_i]$, which satisfies $L_i\leq\lambda_i(A+B)\leq U_i$ by Theorem \ref{thm:weyl-eigenvalue-inequalities-hermitian-sums}(i)--(ii);
\item for structured $B$, simplify the candidates with any known sign or zero-eigenvalue data.\\
If $B\succeq0$, then $\operatorname{spec}(B)\subseteq[0,\infty)$ by Theorem \ref{thm:positive-semidefinite-matrix-equivalences}(ii).\\
If $\operatorname{rank}(B)=r$, then the algebraic multiplicity of $0$ is $n-r$ by Corollary \ref{cor:matrix-rank-nullity}(i) and Theorem \ref{thm:complex-spectral-theorem}.
\end{enumerate}
\end{computationmethod}
